\documentclass[letterpaper]{article}

\usepackage[margin=1.5cm]{geometry}
\usepackage{amsmath,amsfonts}
\usepackage{url}
\usepackage{bm}

\newcommand{\Real}{\mathbb{R}}
\newcommand{\tr}{^\top}
\newcommand{\opt}{^\star}

\setlength{\parskip}{3mm plus1mm minus1mm}
\setlength{\parindent}{0mm}

\title{Problem Set 2}
\author{Foundations of Machine Learning}
\date{}

\begin{document}
\maketitle

\begin{itemize}
\item \textbf{Submission}: Turn in as a \texttt{PDF} and a \textbf{separate source code file} on \url{http://mycourses.unh.edu}. Scanned hand-written notes are encouraged.
\item \textbf{Questions}: See the class website for a link to office hours
\end{itemize}

% Covers: Vectors (ILA 1.1--1.3) and Matrices (ILA 2.1, 2.4, 2.6).

\textbf{Relevant reading}: ILA 1.1--1.3, 2.1, 2.4. Problems marked \textbf{[code]} should be submitted as a program in a language of your choice (Julia or Python recommended). The cited problem numbers are from ILA, 5th edition (Strang, 2016).

\begin{enumerate}
\item ILA, Problem Set 1.1 (vectors and linear combinations): problems 1, 3, 4, 5, 11, 17, 18, 27, 28.
\item ILA, Problem Set 1.2 (lengths and dot products): problems 1, 2, 4, 10, 14, 32.
\item ILA, Problem Set 1.3 (matrices): problems 1, 4.
\item ILA, Problem Set 2.1 (vectors and linear equations): problems 1, 9, 11, 18.
\item \textbf{[code]} Read a CSV data table into a matrix. Compute the dot product of two of its columns, and use a matrix--vector product to normalize every row so that its entries sum to one.
\item \textbf{Time spent.} Report approximately how many hours you spent completing this problem set (including the reading).
\end{enumerate}

\end{document}

%%% Local Variables:
%%% mode: LaTeX
%%% TeX-master: "assignment2"
%%% End:
